\section{Boundary comparisons and separation near the surface}\label{sec:contours}
Put $k_n=\lfloor n^{3/4}\rfloor$ and
\[
 J_n=\{-k_n,\ldots,k_n\}\times\{0\}.
\]
Besides the unrestricted partition function $Z_n$, use $Z_n^+$ and $Z_n^-$, with every spin of $J_n$ fixed respectively to plus and minus. All three retain the physical boundary bonds. Superscripts $+$ and $-$ on these partition functions always mean the line constraint. Define
\begin{equation}\label{eq:F}
 F_n(\eta)=\log Z_n^+(\eta)-\log Z_n^-(\eta).
\end{equation}

\subsection{Concentration of interval comparisons}
Order the $M_n$ exterior bonds cyclically around the square, with the two bonds at a corner in their perimeter order. For a cyclic interval $I$, let $\eta^I$ reverse precisely its exterior spins. The empty and full intervals are allowed. For any one of the three partition functions write
\[
 D_I(\eta)=\log Z(\eta^I)-\log Z(\eta),\qquad m_I=|I|.
\]
Changing one exterior bit changes $\log Z$ by at most $2K$, and hence $D_I$ by at most $4K$.

We shall also need a slightly dependent law on the boundary. Fix an odd integer $b=2L+1$. Partition the first $b\lfloor(2\lfloor n/2\rfloor+1)/b\rfloor$ sites of the middle half $\{(x,n+1):|x|\le\lfloor n/2\rfloor\}$ of the north collar, read from west to east, into consecutive blocks of length $b$. All remaining bits are singletons. In a block the $2L$ noncentral bits are independent fair signs. Call its centre active if all those bits are plus. Conditional on the noncentral bits, an active centre is plus with probability $q$, and an inactive centre is fair. Different blocks and singletons are independent; the latter are fair. Denote this law by $\Pp_q^{(L)}$. At $q=1/2$ it is exactly the fair product law on every bit.

\begin{lemma}[Interval-ratio bounds]\label{lem:ratios}
For product boundary bias $|p-1/2|\le n^{-2/5}$, and also for the block law with $|q-1/2|\le n^{-2/5}$, the following hold. In the product case take $b=1$.
\begin{enumerate}
\item There is an event $H_n$ such that every $D_I$ in the three ensembles is at most $n^{7/10}$ on $H_n$, and
\begin{equation}\label{eq:H}
 \Pp(H_n^c)\le Cn^2\exp\{-n^{2/5}/(288K^2b^2)\}
\end{equation}
for all sufficiently large $n$.
\item If $m_I\le\ell+2$ and $\ell\ge2$, then, for all sufficiently large $n$ uniformly in $I$ and $\ell$,
\begin{equation}\label{eq:localratio}
 \Pp(D_I>t\ell)\le\exp\{-t^2\ell/(64K^2b^2)\}.
\end{equation}
\end{enumerate}
\end{lemma}
\begin{proof}
We give the concentration calculation. If independent coordinates change a function $G$ by at most $d_j$, its Doob martingale increments have conditional ranges at most $d_j$. For a centred variable in an interval of length $d$, the logarithm of its moment generating function has second derivative at most $d^2/4$ (the variance under the tilted law); integration twice gives an upper bound $s^2d^2/8$. Iterating conditional expectations and optimizing Markov's inequality yields
\begin{equation}\label{eq:bd}
 \Pp(G-\E G\ge v)\le
 \exp\left\{-\frac{2v^2}{\sum_j d_j^2}\right\},
 \qquad \Var(G)\le\frac14\sum_jd_j^2.
\end{equation}
The variance bound follows as well by orthogonality of martingale increments.

For product spins, condition on the bits outside $I$. Reversing $I$ replaces bias $p$ by $1-p$. Coordinatewise coupling of the two product distributions changes an expected $|2p-1|m_I$ bits, giving
\[
 |\E(D_I\mid I^c)|\le2K|2p-1|m_I.
\]
The $m_I$ oscillations are at most $4K$, so the upper tail beyond this conditional mean by $v$ is at most $\exp\{-v^2/(8K^2m_I)\}$.

For blocks, condition only on complete groups disjoint from $I$. There are at most $m_I$ intersected groups, including partially intersected endpoint blocks. At fairness their remaining bits are all independent fair, and reversing $I$ preserves this conditional law, so the conditional mean of $D_I$ is zero. Couple a general $q$ to $1/2$, sharing all buffers and centre uniforms. Only a centre in an intersected block can change, with probability at most $|q-1/2|$. A changed centre, even outside $I$ in an endpoint block, changes $D_I$ by at most $4K$. Consequently
\begin{equation}\label{eq:blocktail}
\begin{aligned}
 |\E(D_I\mid\text{disjoint groups})|&\le4Km_I|q-1/2|,\\
 \Pp(D_I-\E(D_I\mid\text{disjoint groups})\ge v\mid\cdots)
 &\le e^{-v^2/(8K^2b^2m_I)}.
\end{aligned}
\end{equation}
A whole group has oscillation at most $4Kb$. This proves the displayed estimate by \eqref{eq:bd}. The same mean upper bound works for the product law.

For $n\ge4$, $M_n\le9n$. The conditional mean is at most $36Kn^{3/5}$, eventually at most $n^{7/10}/2$. Taking $v=n^{7/10}/2$ in \eqref{eq:blocktail} and $m_I\le9n$ gives the exponent in \eqref{eq:H}. There are at most $M_n(M_n+1)$ intervals and three ensembles. Empty intervals contribute zero.

For the local estimate, $m_I\le\ell+2\le2\ell$ and the mean is at most $8Kn^{-2/5}\ell$. Eventually it is at most $t\ell/2$. Use $v=t\ell/2$ in \eqref{eq:blocktail}; this gives \eqref{eq:localratio}. The cutoffs depend on the fixed $K,t,b$, not on $\ell$ or the interval.
\end{proof}

\subsection{The contour injection}
Draw unit square cells centred at the sites of $B_n$. Each disagreeing nearest-neighbour interior bond is crossed by a dual edge. Draw only these interior edges, so the curves end at the perimeter of the union of the cells. Resolve each four-edge contact by a fixed noncrossing pairing, chosen to agree with the interior deformation convention in \cite[Appendix~A.2]{CV}. It can be effected in disjoint small disks about the dual vertices without moving any crossing of a primal bond. In this split geometry each component is a Jordan curve or a simple boundary-to-boundary arc. Its edge walk never repeats an edge or immediately reverses; it may visit a split dual vertex twice.

\begin{lemma}[Flip comparison]\label{lem:flip}
For any specified resolved curve $\gamma$ of $\ell$ interior edges, choose one of its sides. In a line-pinned ensemble choose the side disjoint from $J_n$. Let $I$ be the perimeter interval incident to that side. For a closed curve it is empty on the inside and full on the outside. Then
\begin{equation}\label{eq:flip}
 \mu^\eta(\gamma\text{ occurs})\le e^{-2K\ell}\frac{Z(\eta^I)}{Z(\eta)}.
\end{equation}
For every $\eta\in H_n$ and in each of the three ensembles, the Gibbs probability that some resolved curve of length at least $n^{3/4}$ occurs is at most $Cn^2\rho^{n^{3/4}}/(1-\rho)$ for large $n$; for arcs of length at least $n/4$ the bound is $Cn^2\rho^{n/4}/(1-\rho)$.
\end{lemma}
\begin{proof}
For fixed $\gamma$ and side, flip every interior spin on that side and all incident exterior spins on its perimeter interval. The crossed interior bonds are precisely the $\ell$ edges of $\gamma$. Before the flip they disagree and afterwards they agree. Every other interior bond has both endpoints flipped or both unchanged. Every exterior bond has its two endpoints changed together or unchanged. At a corner the two exterior endpoints are distinct and are handled separately. Thus the new weight is exactly $e^{2K\ell}$ times the old weight. The map is injective, since the same fixed spin subset reverses it; it need not preserve the resolution rule in its image. Summing image weights in the partition function at $\eta^I$ proves \eqref{eq:flip}.

The connected monochromatic line cannot cross a disagreement curve, so the permitted side exists. A loop containing the line requires the outside side and hence the full collar interval; that case is included in $H_n$.

There are $O(n^2)$ possible roots, four initial edge directions and at most three subsequent choices. The traversed edge word fixes its local pairings, so $Cn^2 3^\ell$ overcounts all resolved curves. On $H_n$, $n^{7/10}\le t\ell$ eventually for $\ell\ge n^{3/4}$. The same inequality holds for $\ell\ge n/4$. Sum \eqref{eq:flip} using $3e^{t-2K}=\rho<1$.
\end{proof}

\subsection{Separation at a fixed depth}
For $|x|\le\lfloor n/2\rfloor$, let
\[
 i=(x,n+1),\qquad y=(x,n),\qquad v=(x,n-D).
\]
The site $v$ will measure separation between the two pinned ensembles. Write $\mu_n^{\eta,+J}$ and $\mu_n^{\eta,-J}$ for them.

\begin{lemma}[Separation near the north side]\label{lem:tooth}
For each sign $s$, there are nonnegative functions $R_{n,v}^s$ such that
\begin{equation}\label{eq:error}
 \mu_n^{\eta,sJ}(\sigma_v\ne s)\le B(K)+R_{n,v}^s(\eta).
\end{equation}
For product or fixed-width block laws in Lemma~\ref{lem:ratios}, uniformly over the indicated positions and bias parameters,
\begin{equation}\label{eq:rem}
 \limsup_{n\to\infty}\E R_{n,v}^s\le C_b e^{-c_bD}.
\end{equation}
In particular the ordered plus-probability gap $d_v$ satisfies
\begin{equation}\label{eq:gap}
 d_v\ge g-R_{n,v}^+-R_{n,v}^-.
\end{equation}
\end{lemma}
\begin{proof}
If $v$ disagrees with the pinned line, the parity of disagreements along a path from $v$ to $J_n$ gives a separating curve. A closed separator either encloses $v$ and excludes the line, or encloses the line and excludes $v$. In the former case its interior flip changes no exterior bit, so its cost is $e^{-2K\ell}$ deterministically. A horizontal ray from $v$, which misses dual vertices, meets such an enclosing curve within at most $\ell+1$ possible first crossing edges: a curve of length $\ell$ cannot extend farther in either coordinate from a site it encloses. Orienting the crossing and choosing the subsequent nonbacktracking steps gives the overcount $4(\ell+2)3^{\ell-1}$. The resulting contribution is at most
\begin{equation}\label{eq:Bsum}
 \sum_{\ell\ge4}4(\ell+2)3^{\ell-1}e^{-2K\ell}=B(K).
\end{equation}
This counts odd lengths as well and remains an upper bound.

A loop containing $J_n$ has length at least $2k_n$ and is controlled by the long-curve argument of Lemma~\ref{lem:flip}; its proof applies to any fixed positive multiple of $n^{3/4}$. Open separators of length at least $n/4$ have the second bound there. Their probabilities have expectations tending to zero, by \eqref{eq:H}.

Consider an arc of length $\ell<n/4$. It cannot join opposite sides. Close it by the shorter perimeter route between its endpoints. For endpoints on adjacent sides, that route runs through their common corner. The cap so formed lies in the coordinate bounding rectangle of the path and that perimeter route: any point outside that rectangle can reach infinity along a coordinate ray without crossing the closed curve. Its coordinate span is at most $\ell+O(1)$. Such a cap cannot contain $J_n$, whose distance from the perimeter exceeds $n/4$ for large $n$. To contain the north site $v$ with $|x|\le n/2$ at fixed $D$, a cap of length less than $n/4$ cannot reach a corner or another side. Thus its two endpoints are on the north side, their abscissae are within $\ell+O(1)$ of $x$, and the cap reaches depth $D-O(1)$. With the cell convention above it suffices to sum over $\ell\ge D-2$. The incident collar interval has $m_I\le\ell+2$, since endpoint horizontal displacement is at most the path length, with at most two endpoint bonds added.

There are at most $C(\ell+1)^2$ candidate intervals at this length. Charge each event $\{D_I>t\ell\}$ once. Lemma~\ref{lem:ratios} bounds the expectation of their sum by
\[
 C(\ell+1)^2e^{-t^2\ell/(64K^2b^2)}.
\]
On their complement, root the curve at one of its possible north endpoints and count its edge words. Their total Gibbs probability is at most $C(\ell+1)^3\rho^\ell$, where the extra polynomial factor harmlessly covers endpoint choices. The two tails summed over $\ell\ge\max(2,D-2)$ are bounded by $C_be^{-c_bD}$, after increasing $C_b$ for bounded $D$. In particular the bad-disorder probability is summed over intervals before any curve count; it is not multiplied by $3^\ell$.

Take $R_{n,v}^s$ to be the sum of these bad-interval indicators, the deterministic good-cap tail, and the conditional probabilities of the long separators. This is nonnegative and bounds all contributions apart from \eqref{eq:Bsum}. Its expectation has the bound \eqref{eq:rem}. Finally $d_v=1-\mu_n^{\eta,+J}(\sigma_v=-)-\mu_n^{\eta,-J}(\sigma_v=+)$ gives \eqref{eq:gap}.
\end{proof}
